\chapter{Conclusions}
\indent

The purpose of this work was to understand nitrogen doped chemical vapor deposition of diamond such that optimized conditions could be selected for a desired outcome. In this case, the desired outcome is an ideal substrate for laser writing of nitrogen vacancy (NV) defect arrays which consists of a high quality single crystal diamond plate with a delta doped layer of single substitutional nitrogen. With an ideal substrate, spatial control of NV centers via laser writing will be significantly improved thus improving the spatial resolution of quantum sensors. Determining optimized conditions included understanding the effects of substrate surface preparation on deposition quality, establishing growth quality in a low growth rate regime which expanded the growth operation region of DS4, and a parametric study of reactor parameter effects on nitrogen incorporation. A method to determine nitrogen concentration from FTIR was developed to evaluate deposition quality. A summary of each investigation and concluding remarks on the success of each investigation as well as proposed future direction of each work will be presented here.

\section{Quantification of High Concentration of Substitutional Nitrogen via Fourier Infrared Transform Spectroscopy}
\indent

An accurate, non-destructive method for determining substitutional nitrogen concentration as opposed NV$^-$ was needed to optimized deposition conditions for laser writing substrates. Present methods were either destructive, lacked specificity, or were optimized for low to moderate concentration of nitrogen. The method was developed using FTIR absorption spectra for diamond films grown via CVD with high gas phase concentration of nitrogen. The grown films contained numerous nitrogen-related defects. An accurate, non-destructive method for determining substitutional nitrogen concentration in nitrogen-rich diamond was developed by deconvoluting nitrogen-related IR absorption peaks to determine peak area. Through deconvolution, absorption peaks parameters can be more accurately determined in the presence of strain and other nitrogen-related defects. EPR was used to validate the new method.

\section{Substrate Surface Preparation}
\indent

Several methods for substrate surface preparation were evaluated in Chapter \ref{chSurfacePrep}. First, every substrate was required to have surface roughness of not more than 1 nm after mechanical polishing. Then, an extended hydrogen plasma etch was incrementally administered to a single substrate. Microscale surface features and nanoscale surface roughness were evaluated after each increment. The hydrogen plasma etch preferentially etched polishing damage and dislocations in the substrate. This increased the surface roughness of the sample, after an slight decrease in the first 15 minute etch increment, as the etch time progressed. The same experiment was performed on a second substrate with fewer etch steps. Again, the surface roughness increased as cumulative etch time increased. The increase in surface roughness was a result of increased etch pits with exposed crystalline facets. The removal of amorphous polishing damage and exposure of crystalline facets resulted in low strain grown films. Lastly, a reactive ion etch was performed on the same starting substrate as the first hydrogen etch experiment, however, the grown film was removed and the surface was mechanically polished to not more than 1 nm surface roughness. The reactive ion etch ballistically removed 6.5 $\mu m$ of polishing damaged material from the center of the diamond substrate exposing clean single crystal diamond. The reactive ion etch preferentially etch the polishing lines, but exposed a smooth, low roughness, diamond surface. The reactive ion etch also produced a low strain high quality single crystal diamond film. The extended hydrogen plasma etch was implemented as the surface normalization process due to its ease of access during chemical vapor deposition of diamond, however, reactive ion etching demonstrated equivalent success in surface normalization.

\section{Low Methane Depositions}
\indent

Manufacturing of layered structures through CVD growth potentially requires the ability to slow down or halt growth rates during the deposition for changes in feed gas composition. The reactor used for the work presented in this dissertation was optimized for high growth rates of large single crystal diamond wafers. An experimental deposition series to expand the growth regime of the reactor was completed and high quality SCD growth was validated at the low growth rate conditions. Three samples at 3\%, 2\%, and 1\% methane with no intentional nitrogen doping were grown. Step flow growth, an indicator of high quality growth, was observed in all three depositions. Unepitaxial crystallites formed on the 1\% methane deposition indicating faster growth propagation from nucleation sites than overall vertical growth rate. In other words, the overall vertical plane growth rate did not outpace the growth of adatom nucleation site from early in the deposition process.

In addition to growth morphology, NV background was measured using photoluminescent spectroscopy. Each spectra was normalized to the Raman line allowing for qualitative comparison of intensities. As the methane concentration decreased, the NV concentration increased. This is due to the increased effect of non-negligible atmospheric leak rate into the reactor which increases the ratio of nitrogen to carbon as the methan concentration is decreased, allowing more nitrogen to incorporate in the SCD growth.

The high NV background at reduced growth rates defeats the purpose of utilizing slow growth rates to exchange feed gases and improve resolution of layered structures by minimizing transition layer thickness. The high growth rate reactor is therefore inadequate for layered structure growths due to large residence times and atmospheric nitrogen leak rate. Preliminary investigations into post-processing steps, including low pressure annealing and exposure to a high temperature hydrogen environment, did not reduce NV background of low methane depositions to a usable level for laser writing. 

\section{Effects of Pressure and Offcut on Nitrogen Incorporation}
\indent

A diamond substrate with a sufficient concentration of substitutional nitrogen and low background of NV$^-$ centers is necessary for deterministic laser writing arrays of single NV$^-$ centers. To preferentially incorporate nitrogen as substitutional nitrogen, a parametric study of reactor parameters was conducted to identify deposition conditions to suppress the incorporation of nitrogen as NV$^-$ and favor the incorporation of nitrogen as substitutional nitrogen. Deposition parameters evaluated include substrate temperature, reactor pressure, and substrate offcut. The first experimental series emphasized a correlation in nitrogen incorporation and the confounding variable of surface roughness.

The confounding variable of surface roughness was mitigated through a surface normalization procedure developed in Chapter \ref{chSurfacePrep}. The surface normalization procedure consisted of a mechanical polish of not more than 1 nm surface roughness followed by a 5 hour hydrogen plasma etch immediately prior to deposition. The second experimental series showed no correlation between nitrogen related defect incorporation and surface roughness indicating an effective surface normalization procedure.

Using the second experimental series with normalized surface condition, the effects of reactor pressure and substrate offcut on nitrogen incorporation were investigated. A high nitrogen concentration in the gas phase of the deposition was used to exaggerate the effects of deposition parameter variation on nitrogen incorporation. At high nitrogen concentrations, the sensitivity of reactor pressure and substrate surface orientation was too low to favor the incorporation of nitrogen as substitutional nitrogen and suppress the incorporation of nitrogen as NV$^-$. The NV$^-$ background of the surface normalized deposition series was too high for laser writing. This necessitates the exploration of reactor parameters with higher sensitivity.

The resolution of quantum sensors with NV$^-$ defects can be improved through addressing materials engineering challenges through creation of ideal substrates for laser writing. The manufacturing of these substrates requires precise control of gas phase composition to minimize transition layer thickness of a delta doped region. Further control of nitrogen defect selectivity within the delta doped layer is necessary to reduce NV$^-$ background and incorporation of single substitutional nitrogen.

\section{Proposed Future Work}
\indent
 
The conclusions of this work are that the current CVD reactor, DS4, is not suitable for creating delta doped layers, due to long residence times, high growth rate, and unitentional nitrogen concentration. A different reactor design, such as the NIRIM style reactor, would be better suited for delta doping due to the laminar flow rate of deposition gases. 

Reactor pressure and substrate offcut exhibited low sensitivity with respect to nitrogen defect selection at high gas phase nitrogen concentration. Exploration of other deposition parameters with higher sensitivity in needed. Potential deposition parameters include: low gas phase nitrogen concentration or different nitrogen sources.

Lastly, accurate characterization of nitrogen related defects through FTIR is valuable to non-destructive determination of single substitutional nitrogen concentration. The method developed deconvolutes a complicated absorption band in FTIR using seven Gaussian peaks and was validated using EPR spectroscopy. To ensure the seven nitrogen related peaks used for the fit are present, FTIR data can be taken at low temperatures to reduce the thermal broadening of the peaks allowing each peak to be more easily resolved. For this, a low temperature attachment for the FTIR is necessary.

This continued work would improve the quality of single crystal diamond laser writing substrates with a delta doped substitutional nitrogen layer. The improvement of the laser writing substrate will improve the resolution of NV$^-$ based quantum sensors.